\section{El ecosistema moderno de alineación abierta: más participantes, pero todavía faltan datos}

Las secciones anteriores siguieron los modelos en orden cronológico; esta sección toma abril de 2024 como una instantánea de la clase. Los participantes ya no son solo universidades, organizaciones sin fines de lucro y la comunidad de Hugging Face, sino también numerosas empresas; las recipe de los modelos son más complejas, la evaluación es más rápida y la brecha entre modelos abiertos y cerrados se ha reducido en algunas tareas. Sin embargo, el ponente considera que la limitación más persistente sigue siendo los datos, en particular la preference data humana legal, abierta y diversa.

\subsection{Más modelos y más roles en la generación de datos}

Esta subsección muestra la diversidad del ecosistema moderno. Distintos equipos se especializan en instruction rewriting, preference generation, filtrado, chat data, reward model y frameworks de entrenamiento. La alignment abierta ya no es un trabajo de extremo a extremo de un solo equipo, sino una artifact supply chain.

\lecturefigure{slide-071.jpg}{Diversificación de los roles de generación de datos, entrenamiento de modelos y comunidad en el ecosistema moderno}{official deck, p.~71}

\begin{knowledgebox}{Artifact supply chain}
Un aligned model abierto puede depender de: base weights de terceros, instructions generadas por otro modelo, chat logs públicos, preference data de la comunidad, un trainer de código abierto, un LLM judge cerrado y un leaderboard público. Cada dependencia introduce riesgos de licencia, contaminación, estilo y reproducibilidad, por lo que la model card debe registrar el lineage completo.
\end{knowledgebox}

\subsection{Llama 3: sobre todo una scale story, pero también una post-training signal}

Llama 3 se publicó en las fechas de la clase, y el ponente lo considera «más sobre scaling» que un único breakthrough de alignment. Meta reporta el uso de métodos en múltiples etapas, como IFT, rejection sampling, DPO y PPO; un pipeline tan complejo muestra que la industria está dispuesta a apilar métodos para obtener mejoras pequeñas, y también que todavía no se ha encontrado un algoritmo simple y unificado.

\lecturefigure{slide-072.jpg}{Llama 3: instantánea de la clase sobre el preentrenamiento a gran escala y el post-training en múltiples etapas}{official deck, p.~72}

Conceptualmente, la rejection sampling puede escribirse así: se muestrean $K$ candidatos de la policy
\[
y_1,\ldots,y_K\sim\pi(\cdot\mid x),
\qquad y^{\star}=\arg\max_{y_i} r(x,y_i),
\]
donde $r$ puede ser un reward model, un conjunto de reglas o un judge. Reincorporar $y^{\star}$ al IFT puede mejorar la inicialización, pero también fija en los nuevos datos los sesgos del reward o del judge.

\teachervoice{01:02:44--01:03:22, a Nathan le parece compleja la receta de Llama 3, que combina IFT, rejection sampling, DPO y PPO; conjetura que estas etapas van desplazando gradualmente las capacidades y sirven de inicialización para la etapa siguiente, y que en el futuro podrían destilarse en un algoritmo más simple.}

\subsection{Current directions: la alignment se está convirtiendo en un data problem}

Esta subsección pasa de los modelos a las preguntas de investigación. La comunidad abierta cuenta con pocos preference datasets de uso extendido, como Anthropic HH, UltraFeedback y Nectar; su reutilización repetida provoca sobreajuste, homogeneización y contamination. En el futuro se necesitarán más datos de interacción reales con personas, multilingües, de múltiples dominios y de largo plazo.

\lecturefigure{slide-073.jpg}{Current directions: la frontera de investigación en la alineación de modelos de lenguaje abiertos}{official deck, p.~73}

\lecturefigure{slide-074.jpg}{La brecha de capacidades entre aligned models abiertos y cerrados varía según la tarea}{official deck, p.~74}

\begin{warningbox}{«Open catches closed» debe especificar la tarea y la fecha}
Que algunos modelos abiertos se acerquen a los cerrados en benchmarks públicos no significa que los hayan alcanzado en uso de herramientas, contexto largo, confiabilidad, seguridad, latencia y capacidad multilingüe. La figura es un snapshot de abril de 2024; las conclusiones cambian al actualizarse las versiones de los modelos y los judges. Toda comparación debe indicar el exact model, la date y el evaluation protocol.
\end{warningbox}

\lecturefigure{slide-075.jpg}{Primera dirección actual: la preference data es muy insuficiente y la calidad de los datos limita los experimentos}{official deck, p.~75}

\begin{importantbox}{Matriz de cobertura de un preference dataset}
Los datos de alta calidad deben cubrir, como mínimo: dominio de la tarea, idioma, grupo de usuarios, nivel de riesgo, longitud de la respuesta, con herramientas o sin ellas, estado de múltiples turnos y disagreement. Aumentar solo el número de muestras de una misma synthetic template reduce la varianza, pero no amplía los límites del comportamiento.
\end{importantbox}

\teachervoice{01:00:00--01:03:48, el ponente señala la escasez de preference data como el cuello de botella central actual, y advierte que la synthetic data suele ser repetitiva y de distribución inestable, por lo que su aporte a la generalización de la alignment es limitado.}

\subsection{Dónde ocurre la alineación abierta}

La última diapositiva de contenido enumera participantes como AI2, HuggingFaceH4, Berkeley, LMSYS y NousResearch. La lista no pretende ser una clasificación completa, sino mostrar que el progreso abierto está distribuido en varios nodos: modelos, datos, trainer, evaluation y gestión de la comunidad; un ecosistema sano requiere que estos artifacts puedan combinarse entre sí y auditarse.

\lecturefigure{slide-076.jpg}{Principales comunidades, instituciones y tipos de artifact de la alignment abierta}{official deck, p.~76}

\begin{knowledgebox}{Lectura de la figura: la división del trabajo detrás de la lista de instituciones}
AI2 se centra en Tulu/OLMo y en la publicación de datos; HuggingFaceH4 reproduce recipes con rapidez; LMSYS aporta Arena y datos de conversaciones reales; Berkeley y los equipos de la comunidad exploran PPO y la evaluación; otros equipos generan instructions o preference data. La ventaja del ecosistema abierto es la exploración en paralelo; su debilidad es que la provenance, las licencias y los criterios de evaluación no están unificados.
\end{knowledgebox}

\teachervoice{La recomendación de investigación de 01:04:49--01:05:30 es: nadie puede seguir todos los avances, así que conviene seguir desarrollando habilidades y construir lo que uno considere importante. Para el ecosistema abierto, esto significa que aportar datos, evaluaciones o herramientas de alta calidad es igualmente investigación central, sin necesidad de perseguir solo la escala de los modelos.}

\subsection{Synthetic data y su relación con los datos humanos}

Los datos sintéticos permiten sortear parte de los problemas de derechos de autor, privacidad y costo de anotación, y también generar tokens a mayor escala; sin embargo, los sesgos del teacher model, los patrones repetitivos y la falta de intent de usuarios reales limitan la generalización. Una dirección más razonable es un pipeline mixto de seed humana, ampliación con modelos, filtrado automático y auditoría manual.

Sea $p_h$ la distribución de los datos reales y $p_s$ la distribución sintética; la distribución de entrenamiento mixta puede escribirse como
\[
p_{\lambda}=\lambda p_h+(1-\lambda)p_s,
\]
donde $\lambda$ controla la proporción de datos humanos. El $\lambda$ óptimo no es una constante fija: cuando cambian la calidad de los datos, el dominio de la tarea o la teacher strength, debe volver a elegirse mediante una held-out human evaluation.

\teachervoice{En la Q\&A, de 01:14:38--01:15:29, Nathan considera que la synthetic data es una de las pocas vías que permiten seguir generando más tokens, y cree que los modelos seguirán participando en el entrenamiento de otros modelos; pero aclara que esta dirección todavía está en una etapa muy temprana.}

\subsection{Resumen de la sección}

El protagonista del ecosistema moderno pasó de ser un modelo único a una artifact network. Llama 3 muestra la combinación de una base a gran escala con post-training en múltiples etapas; la brecha entre modelos abiertos y cerrados varía según la tarea; la preference data de alta calidad y la evaluación confiable siguen siendo recursos escasos. El progreso futuro podría venir de algoritmos unificados más simples, o bien de una mejor infraestructura de datos y de retroalimentación.
